\subsection{USB 串口与 reboot 系统调用}\label{sec:serialreboot}

机械臂的运动控制器并不挂在开发板的板载串口上，而是经一颗 CP210x USB 转串口桥接芯片接入，
在系统里表现为 \texttt{/dev/ttyUSB0} 这样一个普通的终端设备。机器人通过这个节点读取机械臂
的当前状态，并把抓取与归位指令写回去（见 \ref{sec:arm}）。USB 转串口桥的作用是把一段异步
串行的 UART 字节流封装进 USB 的批量传输，主机一侧并不存在真正的 UART 控制器，波特率、数据位、
停止位这些线路参数都不是去写某个串口寄存器，而是经 USB 的控制传输下发给桥接芯片，由它在另一侧
还原成对应的电平时序。因此驱动要做的事分成两层，一层是把芯片配置成所需的线路模式，另一层是把
应用读写的字节与芯片的批量端点对接，并让这一切套进内核既有的终端设备模型。

\paragraph{从描述符识别芯片并定位数据端点}
一个 USB 设备插入后，主机先取回它的设备描述符与配置描述符。本队的可复用驱动核心是一个
\texttt{\#![no\_std]} 的库，它直接在这份原始描述符字节流上工作。识别 CP210x 的依据是设备描述符里的
厂商号 \texttt{0x10c4} 与产品号 \texttt{0xea60}，二者同时匹配才认定。随后它遍历配置描述符下的
各个接口，挑出那个厂商自定义类（接口类 \texttt{0xff}、子类与协议均为零、且为零号备用设置）的数据
接口，再在该接口下按方向各取一个批量端点，得到一进一出的端点地址对。描述符的解析全程做长度与类型
校验，遇到越界或长度为零的项即停止，避免在一份畸形的描述符上读出错位的字段。这一层只产出一个
\texttt{UsbSerialPort}，记录接口号与批量端点地址，不触碰任何具体的传输提交方式。

\paragraph{按厂商约定配置串口参数}
CP210x 的所有配置都走 USB 控制传输里的厂商请求，请求类型固定为厂商方向接口（\texttt{0x40 |
0x01}），请求号则对应芯片手册里的各项操作。打开一条串口时，驱动依次发出启用 UART
（\texttt{IFC\_ENABLE}，值 \texttt{0x0001}）、设置波特率（\texttt{SET\_BAUDRATE}，把波特率以小端
四字节作为数据载荷）、设置线路控制为 8 数据位（\texttt{SET\_LINE\_CTL}，值 \texttt{0x0800}）、
拉起 DTR 与 RTS 并置其写使能位（\texttt{SET\_MHS}），最后下发一段全零的流控配置
（\texttt{SET\_FLOW}）以关闭硬件流控。这套时序与 Linux 上的 cp210x 驱动一致，目的就是把芯片置于
一条不带流控的裸 8N1 链路，让上层只看到透明的字节流。波特率不是写死的常量，而是从终端的
\texttt{termios} 里取，应用通过标准的终端接口改波特率时，驱动会重新下发一次 \texttt{SET\_BAUDRATE}，
机械臂控制器所需的速率因此由用户态决定而非驱动假定。

\paragraph{把字节流接进终端设备}
内核侧的接入复用了已有的 \texttt{Tty} 框架，与板载串口共用同一套行规程与终端语义，区别只在底层
的收发后端。\texttt{/dev/ttyUSB0} 是一个静态存在的设备节点，打开时才惰性地去占用第一个被识别到的
USB 串口适配器，占用一个 usbfs 会话并认领对应接口，再跑前述的初始化时序，使得终端可以先于适配器
插入就存在。会话建立后，一个接收工作线程持续在批量输入端点上做读入，把收到的字节推进接收队列并唤醒
等待的读者；发送侧则把应用写入与终端回显的字节并入同一条发送队列，由发送路径在批量输出端点上排空，
两路输出共用一把输出锁以保证字节顺序不被打乱。当前的 xHCI 后端尚未提供基于 Stop Endpoint 的传输
取消能力，因此一条仍在阻塞读的接收传输不会被强行打断，会话的拆除被推迟到接收线程观察到一次真实的
USB 完成或设备错误时再收尾，这是当前后端能力下的稳妥处理，而非可随时中断的理想形态。\textbf{该驱动以
合并请求 \href{https://github.com/rcore-os/tgoskits/pull/1378}{\#1378} 提交至上游并已合并，host
侧单元测试覆盖了厂商号匹配、非 CP210x 设备的拒识、批量端点配对以及控制时序的逐条核对，并在
OrangePi 5 Plus 上完成了实机验证。}

\paragraph{一条命令切换整套系统}
开发期间只有一块远端开发板，而验证又要求在同一块板上反复切换 Linux 与 Starry，跑同一个二进制做
对照。若每次切换都靠人工断电再上电，节奏会被这一步拖住。为此本队实现了 \texttt{reboot} 系统调用，
切换从一次手动电源循环变成一条命令，迭代因此明显加快。这个调用同样服务于面向人工智能开发的那套
自动化回路（见 \ref{sec:ai}），让镜像替换后的重启不再需要有人守在板边。

实现遵循 Linux 的 \texttt{reboot} 语义。入口先校验调用者持有 \texttt{CAP\_SYS\_BOOT} 能力，再核对两个
魔数，第一魔数须为 \texttt{0xfee1dead}，第二魔数须落在 Linux 认可的几个取值之内，任一不符即以
\texttt{EINVAL} 拒绝，避免误触的写入意外重启整机。命令字区分几种行为，\texttt{RESTART} 与
\texttt{RESTART2} 触发重启，\texttt{HALT} 与 \texttt{POWER\_OFF} 触发关机，\texttt{CAD\_ON} 与
\texttt{CAD\_OFF} 则直接返回成功而不做实际动作。重启与关机之前都会先调用一次文件系统的下刷与卸载，
让缓存中的数据落盘，再把控制交给平台电源层。\texttt{RESTART2} 携带的重启原因字符串当前被忽略，
\texttt{CAD\_ON}／\texttt{CAD\_OFF} 也尚未真正改变组合键的行为，这两点已如实记录为已知限制。

\paragraph{重启最终如何作用于硬件}
平台电源层把重启统一抽象为一个 \texttt{reset()} 操作，具体动作随体系结构而异，在本机所用的 aarch64
上走 PSCI（电源状态协调接口，由更高特权级的固件提供的标准电源服务）。系统启动时，电源模块从设备树
里找到 \texttt{arm,psci-1.0}（含 \texttt{0.2} 与无版本号的回退）节点，读出其 \texttt{method} 属性以
确定调用固件的方式是 \texttt{smc} 还是 \texttt{hvc}，并把它记下。真正重启时，据这一方式发起 PSCI 的
\texttt{SYSTEM\_RESET}，由固件接管整机复位；调用经 \texttt{smccc} 库封装，若固件返回错误则打印诊断
并停在等待中断的空转循环里。关机走的是同一条路径下的 \texttt{SYSTEM\_OFF}。\textbf{该系统调用以合并请求
\href{https://github.com/rcore-os/tgoskits/pull/1358}{\#1358} 提交至上游并已合并。}
